\chapter{Lambda Calcolo}
Ideato da Alonzo Church nel 1931, il $\lambda$-calcolo nasce per esplicitare il \textit{processo meccanico di calcolo} (sintassi) di una funzione, in contrasto con la visione matematica classica che vede la funzione solo come una relazione statica tra insiemi (semantica).

\section{Sintassi e Computazione}
La sintassi del $\lambda$-calcolo è minimalista ed è composta da tre regole:
\begin{itemize}
    \item \textbf{variabile:} $x, y, z \dots$
    \item \textbf{astrazione:} $\lambda x. M$ (corrisponde alla definizione di una funzione con parametro $x$ e corpo $M$)
    \item \textbf{applicazione:} $(M N)$ (corrisponde all'applicazione della funzione $M$ all'argomento $N$)
\end{itemize}
\vspace{0.5em}
in BNF, quindi, la grammatica del lambda-calcolo è quindi:
$$T ::= V \mid \lambda V.T \mid (T \ T)$$

\begin{gframe}{Regole di associazione}
    \begin{itemize}
        \item \textbf{l'applicazione associa a sinistra:} $F N_1 N_2 \dots N_n$ equivale a $(\dots (F N_1) \dots N_n)$
        \item \textbf{l'astrazione associa a destra:} $\lambda x_1 \dots x_n. M$ equivale a $\lambda x_1. (\lambda x_2. (\dots (\lambda x_n. M)\dots))$
    \end{itemize}
\end{gframe}

La computazione avviene tramite la \textbf{$\beta$-riduzione} (beta-regola): un termine della forma $(\lambda x.M) N$ si dice \textit{redex} e si riduce sostituendo tutte le occorrenze di $x$ in $M$ con il termine $N$.
Formalmente:
$$(\lambda x.M) N \rightarrow_\beta M[N/x]$$

Un termine che non contiene più $\beta$-redex è detto in \textbf{forma normale} e rappresenta un valore finale (il risultato della computazione).

\section{Variabili e Combinatori}
\begin{itemize}
    \item \textbf{Variabili libere ($FV$) e legate:} in $\lambda x. M$, la variabile $x$ è \textit{legata} (bound) all'astrazione. Qualsiasi variabile non legata è \textit{libera} (free).
    \item \textbf{$\alpha$-conversione (alfa-regola):} per evitare la cattura accidentale di variabili libere durante la sostituzione, si possono rinominare le variabili legate (es. $\lambda x.x \equiv \lambda y.y$). 
        \begin{itemize}
            \item è buona pratica seguire la ``convenzione di Barendregt'', che stabilisce che, all'interno di un termine lambda o di un insieme di termini, tutte le variabili legate devono avere nomi distinti tra loro e diversi da tutte le variabili libere presenti nel termine stesso
        \end{itemize} 
    \item  I termini tali che ($FV(M) = \emptyset$) ($\lambda$-termini chiusi), ovvero senza variabili libere, sono chiamati \textbf{combinatori}
\end{itemize}

\subsection{Combinatori noti e strutture di controllo}
Il $\lambda$-calcolo puro non ha tipi di dato predefiniti: tutto (booleani, numeri, if-then-else) deve essere codificato tramite funzioni.

\begin{itemize}
    \item \textbf{Identità ($I$):} $I \equiv \lambda x.x$
    \item \textbf{Cancellatori (booleani):}
        \subitem \textbf{True ($K$):} $K \equiv \lambda x y. x$ (prende due argomenti e restituisce il primo)
        \subitem \textbf{False ($O$):} $O \equiv \lambda x y. y$ (prende due argomenti e restituisce il secondo)
    \item \textbf{If-Then-Else:} grazie alla definizione dei booleani, la struttura condizionale è codificata come un'applicazione: $if\ x\ then\ M\ else\ N \equiv x M N$. Se $x$ è True ($K$) selezionerà $M$, se è False ($O$) selezionerà $N$.
    \item \textbf{Compositore con condivisione ($S$):} $S \equiv \lambda x y z. x z (y z)$ - distribuisce l'argomento $z$ sia alla funzione $x$ che alla funzione $y$
\subitem $S K K \approx I$: applicando un termine generico $w$, si ha $(S K K) w \rightarrow K w (K w)$. Poiché $K w$ restituisce la funzione costante che ignora il secondo argomento, $K w (K w) \rightarrow w$. Di conseguenza $(S K K) w \rightarrow w$, che equivale a $I w$.
    \item \textbf{Duplicatori e divergenza:} 
        \subitem $\omega \equiv \lambda x.xx$ (auto-applicazione)
        \subitem $\Omega \equiv (\lambda x.xx) (\lambda x.xx)$, ovvero $\omega \omega \rightarrow \omega \omega \rightarrow \dots$ rappresenta la \textbf{funzione indefinita}, il loop infinito (non tipabile)
\end{itemize}

\subsubsection{Numeri di Church e iteratori}
I \textbf{Numerali di Church} rappresentano i numeri naturali come funzioni di ordine superiore che \textit{iterano} un'operazione. Il numero $n$ è una funzione che prende una funzione (es. successore) $s$ e un valore di partenza (es. zero) $z$, e applica $s$ ad $z$ esattamente $n$ volte
$$\underline{n} \equiv \lambda s z. s(s(\dots(sz)\dots)) \quad [n\text{ applicazioni}]$$
\begin{itemize}
    \item $\underline{0} \equiv \lambda s z. z$ (equivale a False ($K$)) 
    \item $\underline{1} \equiv \lambda s z. s z$
    \item $\underline{2} \equiv \lambda s z. s(s(z))$
    \item $\underline{3} \equiv \lambda s z. s(s(s(z)))$
\end{itemize}

Per calcolare il \textbf{successore} ($+1$) di un numero di Church, basta creare una funzione che applica $s$ una volta in più: $succ \equiv \lambda x s z. s(x s z)$ (o anche $\lambda x s z. xs (s z)$)

I numerali di Church non sono quindi solo dati, ma \textbf{iteratori intrinseci}.

\section{Ricorsione e Combinatori di Punto Fisso}
Per definire funzioni ricorsive complesse il $\lambda$-calcolo necessita di un principio di ricorsione generale. 

\begin{thmframe}{Combinatore di Punto Fisso}{}
 Nel $\lambda$-calcolo esiste un termine $Y$ (detto \textbf{Combinatore di Punto Fisso}) tale che, per ogni termine $M$, si ha $Y M \rightarrow M (Y M)$. 
 
 Il combinatore permette di trovare quindi il punto fisso di una funzione (quel valore $x$ tale che $f(x) = x$)
\end{thmframe}

Due combinatori di punto fisso sono:
\begin{itemize}
    \item \textbf{Combinatore di Turing ($Y_T$):} definito come $Y_T = \theta \theta$ dove $\theta \equiv \lambda x y. y(x x y)$.
\begin{align*}
Y_T M &\equiv (\theta \theta) M \\
&\equiv ((\lambda x y. y(x x y)) \theta) M \\
&\rightarrow_\beta (\lambda y. y(\theta \theta y)) M \\
&\rightarrow_\beta M (\theta \theta M) \\
&\equiv M (Y_T M)
\end{align*}

    \item \textbf{Combinatore di Curry ($Y_C$):} definito come $Y_C \equiv \lambda f. (\lambda x.f(xx))(\lambda x.f(xx))$. Se applicato all'identità $I$, questo combinatore riduce a $\Omega$ (divergenza).
    \begin{align*}
    Y_C M &= (\lambda f. (\lambda x. f(x x)) (\lambda x. f(x x))) M \\
    &\rightarrow_\beta (\lambda x. M(x x)) (\lambda x. M(x x)) \\
    &\rightarrow_\beta M ( (\lambda x. M(x x)) (\lambda x. M(x x)) ) \\
    &\equiv M (Y_C M)
    \end{align*}
    {\small \color{gray}(l'ultima equivalenza $\equiv$ vale perché $(\lambda x. M(x x)) (\lambda x. M(x x))$ è la forma ridotta di $Y_C M$)}

\end{itemize}

\begin{thmframe}{Tesi di Church-Turing}{}
 Ogni funzione calcolabile è $\lambda$-definibile.
\end{thmframe}

